\scanpage{27}
\mspage{13}

De façon précise, \inserted{la donnée} des \inserted{$M_0$-torseurs} $M_i$ et \inserted{de} l'isom can. \unsure{(6)} revient à la donnée du torseur $P_M$ sous $V_J(M_0)$, avec
\[
M_i\simeq P_M\wedge_{M_0^J}(M_0,\operatorname{pr}_i)\qquad\text{pour }i\in J \tag*{(9)}
\]
où $\operatorname{pr}_i:M_0^J\mathrel{\text{---}}M_0$ la projection d'indice $i\in J$.

Notons qu'on a des iso. canoniques (valables pour tt vectoriel de rang $2$ sur $\mathbb F_2$)
\[
\begin{aligned}
\operatorname{End}(M_0)&\simeq V_J(M_0)\xrightarrow{\sim}M_0\otimes M_0\simeq M_0\oplus\mathbb F_2^{\underline\omega_{M_0}}\\
\underline\omega_{M_0}&\simeq\underline\omega(\underset{J}{\underbrace{M_0^*}})\quad\text{ens des orientations de }J.
\end{aligned}\tag*{(10)}
\]

% slanted
\margin{\uline{NB} Pour tt vectoriel $W$ sur $\mathbb F_2$ on a (11) $\operatorname{Hom}_{\mathbb F_2}(M_0,W)\simeq V_J(\unsure{W|V})$ Car $M_0\simeq\mathbb F_2^J/\operatorname{diag}$ \deleted{est} $\longrightarrow$}

\[
\left\{\begin{aligned}
\operatorname{End}(M_0)&\xrightarrow{\sim}V_J(M_0)\\
u&\longmapsto(u(x_i))_{i\in J}
\end{aligned}\right. \tag*{(12)}
\]
\[
\operatorname{End}(M_0)\simeq M_0^\vee\otimes M_0\simeq M_0\otimes M_0
\qquad\substack{\text{car }M_0^\vee\simeq M_0\\\text{(valable pour tt}\\\text{vect. de rang }2\text{ sur }\mathbb F_2)}
\]

\diagram{\begin{tikzcd}[ampersand replacement=\&,column sep=small,row sep=0pt]
(13)\& M_0\arrow[r,hook]\&\operatorname{End}(M_0)\\
\&\rotatebox{90}{$\subset$}\&\rotatebox{90}{$\subset$}\\
\&M_0^*\arrow[r,leftrightarrow]\&\operatorname{Aut}(M_0)\simeq\unsure{\mathfrak S|G}_J\\
\&x_i\arrow[r,mapsto]\&\substack{\deleted{\text{transposition}}\ \inserted{\sigma_i}\\\text{réflexion de }M_0\text{ qui fixe }x_i\\
\text{(donc }\deleted{\unsure{\text{qui échange}}}\ \inserted{\text{transposition}}\text{ de }x_j\text{ et }x_k)}
\end{tikzcd}}

[\uline{NB} Comme $\sum_{i\in J}\sigma_i=0$, cela correspond à \unsure{hom}
\[
M_0\longrightarrow\operatorname{End}(M_0)\ ]
\]

\diagram{\begin{tikzcd}[ampersand replacement=\&,column sep=small,row sep=0pt]
(14)\&\mathbb F_2^{\underline\omega_{M_0}}\arrow[r,hook]\&\operatorname{End}(M_0)\\
\&\rotatebox{90}{$\subset$}\&\rotatebox{90}{$\subset$}\\
\&\underline\omega_{M_0}\arrow[r,hook,"\mathrm{can}"]\&\operatorname{Aut}(M_0)\simeq\mathfrak S_J\\
\&\Vert\&\\
\&\underline\omega_J\simeq\mathfrak S_J^+\smallsetminus\{\operatorname{id}_J\}\&
\end{tikzcd}}

\uline{NB} \deleted{\unsure{Comme}} \inserted{On a} \inserted{si $\rho\in\underline\omega$, $\rho'\overset{\mathrm{déf}}=\rho^{-1}$}
\[
\rho+\rho'\deleted{\unsure{=|+}}\,+\unsure{\deleted{1}|1}\operatorname{id}=0
\]
i.e. $\rho+\rho'=\operatorname{id}_{M_0}$

i.e. l'hom $\mathbb F_{\unsure{2|1}}^{\underline\omega}\mathrel{\unsure{\lhook\joinrel\relbar|\hookrightarrow}}\operatorname{End}(M_0)$ applique $1$ sur $1$.

Si $i,j\in J$, $i\ne j$, alors $\sigma_i,\sigma_j,\rho,\rho'\in\operatorname{End}(M_0)$ sont lin. indépendants sur $\mathbb F_2$ donc forment une base, \deleted{et} alors que \unsure{$x_i,x_j$|$\sigma_i,\sigma_j$|\deleted{$x_i,x_j$}} en forment une de $M_0$, et $\rho,\rho'$ de $\mathbb F_2^{\underline\omega}$ \ldots

% archive sheet number: 26

